\documentclass[10pt,a4paper]{amsart}
\usepackage{amsmath,amssymb,amsthm,graphicx}
\usepackage{verbatim}
\newenvironment{reference}{\comment}{\endcomment}
%\newenvironment{reference}{}{}
\newenvironment{slogan}{\comment}{\endcomment}
\newenvironment{history}{\comment}{\endcomment}

% For commutative diagrams we use Xy-pic
\usepackage[all]{xy}

% We use 2cell for 2-commutative diagrams.
\xyoption{2cell}
% Theorem environments.
%
\theoremstyle{plain}
\newtheorem{theorem}[subsection]{Theorem}
\newtheorem{proposition}[subsection]{Proposition}
\newtheorem{lemma}[subsection]{Lemma}

\theoremstyle{definition}
\newtheorem{definition}[subsection]{Definition}
\newtheorem{example}[subsection]{Example}
\newtheorem{exercise}[subsection]{Exercise}
\newtheorem{situation}[subsection]{Situation}

\theoremstyle{remark}
\newtheorem{remark}[subsection]{Remark}
\newtheorem{remarks}[subsection]{Remarks}

\numberwithin{equation}{subsection}

% Macros
%
\def\lim{\mathop{\mathrm{lim}}\nolimits}
\def\colim{\mathop{\mathrm{colim}}\nolimits}
\def\Spec{\mathop{\mathrm{Spec}}}
\def\Hom{\mathop{\mathrm{Hom}}\nolimits}
\def\Ext{\mathop{\mathrm{Ext}}\nolimits}
\def\SheafHom{\mathop{\mathcal{H}\!\mathit{om}}\nolimits}
\def\SheafExt{\mathop{\mathcal{E}\!\mathit{xt}}\nolimits}
\def\Sch{\mathit{Sch}}
\def\Mor{\mathop{\mathrm{Mor}}\nolimits}
\def\Ob{\mathop{\mathrm{Ob}}\nolimits}
\def\Sh{\mathop{\mathit{Sh}}\nolimits}
\def\NL{\mathop{N\!L}\nolimits}
\def\CH{\mathop{\mathrm{CH}}\nolimits}
\def\proetale{{pro\text{-}\acute{e}tale}}
\def\etale{{\acute{e}tale}}
\def\QCoh{\mathit{QCoh}}
\def\Ker{\mathop{\mathrm{Ker}}}
\def\Im{\mathop{\mathrm{Im}}}
\def\Coker{\mathop{\mathrm{Coker}}}
\def\Coim{\mathop{\mathrm{Coim}}}

% Boxtimes
%
\DeclareMathSymbol{\boxtimes}{\mathbin}{AMSa}{"02}

%
% Macros for moduli stacks/spaces
%
\def\QCohstack{\mathcal{QC}\!\mathit{oh}}
\def\Cohstack{\mathcal{C}\!\mathit{oh}}
\def\Spacesstack{\mathcal{S}\!\mathit{paces}}
\def\Quotfunctor{\mathrm{Quot}}
\def\Hilbfunctor{\mathrm{Hilb}}
\def\Curvesstack{\mathcal{C}\!\mathit{urves}}
\def\Polarizedstack{\mathcal{P}\!\mathit{olarized}}
\def\Complexesstack{\mathcal{C}\!\mathit{omplexes}}
% \Pic is the operator that assigns to X its picard group, usage \Pic(X)
% \Picardstack_{X/B} denotes the Picard stack of X over B
% \Picardfunctor_{X/B} denotes the Picard functor of X over B
\def\Pic{\mathop{\mathrm{Pic}}\nolimits}
\def\Picardstack{\mathcal{P}\!\mathit{ic}}
\def\Picardfunctor{\mathrm{Pic}}
\def\Deformationcategory{\mathcal{D}\!\mathit{ef}}
\usepackage{cleveref}
\crefname{originallemma}{original lemma}{original lemmas}
\Crefname{originallemma}{Original lemma}{Original lemmas}
\makeatletter
\newlabel{topology-lemma-constructible-hausdorff-quasi-compact}{{0901}{source}{}{sourcepointer-1}{}}
\newlabel{topology-lemma-constructible-hausdorff-quasi-compact@cref}{{[originallemma][0][]0901}{[1][0][]source}{}{}{}}
\makeatother
\crefname{originallemma}{original lemma}{original lemmas}
\Crefname{originallemma}{Original lemma}{Original lemmas}
\makeatletter
\newlabel{properties-lemma-quasi-compact-quasi-separated-spectral}{{094L}{source}{}{sourcepointer-2}{}}
\newlabel{properties-lemma-quasi-compact-quasi-separated-spectral@cref}{{[originallemma][0][]094L}{[1][0][]source}{}{}{}}
\makeatother
\crefname{originallemma}{original lemma}{original lemmas}
\Crefname{originallemma}{Original lemma}{Original lemmas}
\makeatletter
\newlabel{lemma-constructible-quasi-compact-quasi-separated}{{095E}{source}{}{sourcepointer-3}{}}
\newlabel{lemma-constructible-quasi-compact-quasi-separated@cref}{{[originallemma][0][]095E}{[1][0][]source}{}{}{}}
\makeatother
\crefname{originallemma}{original lemma}{original lemmas}
\Crefname{originallemma}{Original lemma}{Original lemmas}
\makeatletter
\newlabel{topology-lemma-constructible-in-constructible}{{09YJ}{source}{}{sourcepointer-4}{}}
\newlabel{topology-lemma-constructible-in-constructible@cref}{{[originallemma][0][]09YJ}{[1][0][]source}{}{}{}}
\makeatother
\crefname{originallemma}{original lemma}{original lemmas}
\Crefname{originallemma}{Original lemma}{Original lemmas}
\makeatletter
\newlabel{topology-lemma-constructible-partition-refined-by-stratification}{{09Y4}{source}{}{sourcepointer-5}{}}
\newlabel{topology-lemma-constructible-partition-refined-by-stratification@cref}{{[originallemma][0][]09Y4}{[1][0][]source}{}{}{}}
\makeatother
\title{Constructibility is local for a covering by constructible locally closed subschemes}
\author{The Stacks Project Authors}
\date{stacks-project:etale-cohomology.tex:local-snapshot-20261011; GNU FDL 1.2 or later}
\begin{document}
\maketitle
\noindent\textit{Editorial note.} All delivered mathematical text is the author's own, copied from the pinned Stacks Project chapter file etale-cohomology.tex; no statement or proof step was rewritten, reordered or removed.\par\medskip
\noindent\textit{Editorial note.} The proof is the author's complete argument; the sole parenthetical about the details of a piece being quasi-compact and quasi-separated concerns a standard property of constructible subsets and the surrounding argument is written out with all its inputs.\par\medskip
\noindent\textit{Editorial note.} Every labelled reference occurring inside the counted spans is carried in assembly.external\_references as a non-load-bearing navigation pointer (same-chapter ones at their definition site, cross-chapter ones at the invocation inside the counted span), so no load-bearing mathematics is replaced by a stub. No bibliography citation occurs inside the counted spans.\par\medskip
\noindent\textit{Editorial note.} The checked chapter list for same-name collisions is the set of stacks chapters staged in this repository: chow, cohomology, divisors, etale-cohomology, flat, more-algebra, more-morphisms, morphisms, spaces-more-morphisms, varieties together with this card's chapter etale-cohomology.\par\medskip
\noindent\textit{Editorial note.} The preamble is the chapter's own native preamble excerpt already used for this source: the Xy-pic setup (native preamble lines 14-18) and the native theorem declarations with the mathematical macros (native preamble lines 154-217).\par\medskip
\noindent\textit{Allowed standard prerequisites (editorial).}\begin{itemize}
\item Characterization of constructible sheaves by finite partitions. Same chapter, Lemma lemma-constructible-quasi-compact-quasi-separated. Turns constructibility into finite local constancy on the pieces of a partition.
\item Constructible sets form a spectral, quasi-compact topology. Topology, Lemma topology-lemma-constructible-hausdorff-quasi-compact together with Properties, Lemma properties-lemma-quasi-compact-quasi-separated-spectral. Lets the author assume the covering is finite.
\item Constructibility of a constructible set inside a constructible set, and refinement by a stratification. Topology, Lemmas topology-lemma-constructible-in-constructible and topology-lemma-constructible-partition-refined-by-stratification. Produces the common refinement partition of the whole scheme.
\end{itemize}\medskip
\ifdefined\hypertarget\hypertarget{sourcepointer-1}{}\fi
\noindent\textit{Source reference.} Original lemma 0901: sources/stacks-etale-cohomology/pinned-original-source.tex, source lines 11771--11771. Navigation pointer invoked at proof line 11771; the counted span contains the pointer only and the cited result is not reproduced inside it, so the pointer is registered as a non-load-bearing navigation stub in the established form for Stacks material. No bibliography citation is involved.\par\medskip
\ifdefined\hypertarget\hypertarget{sourcepointer-2}{}\fi
\noindent\textit{Source reference.} Original lemma 094L: sources/stacks-etale-cohomology/pinned-original-source.tex, source lines 11774--11774. Navigation pointer invoked at proof line 11774; the counted span contains the pointer only and the cited result is not reproduced inside it, so the pointer is registered as a non-load-bearing navigation stub in the established form for Stacks material. No bibliography citation is involved.\par\medskip
\ifdefined\hypertarget\hypertarget{sourcepointer-3}{}\fi
\noindent\textit{Source reference.} Original lemma 095E: sources/stacks-etale-cohomology/pinned-original-source.tex, source lines 11717--11717. Navigation pointer invoked at proof line 11780; the counted span contains the pointer only and the cited result is not reproduced inside it, so the pointer is registered as a non-load-bearing navigation stub in the established form for Stacks material. No bibliography citation is involved.\par\medskip
\ifdefined\hypertarget\hypertarget{sourcepointer-4}{}\fi
\noindent\textit{Source reference.} Original lemma 09YJ: sources/stacks-etale-cohomology/pinned-original-source.tex, source lines 11781--11781. Navigation pointer invoked at proof line 11781; the counted span contains the pointer only and the cited result is not reproduced inside it, so the pointer is registered as a non-load-bearing navigation stub in the established form for Stacks material. No bibliography citation is involved.\par\medskip
\ifdefined\hypertarget\hypertarget{sourcepointer-5}{}\fi
\noindent\textit{Source reference.} Original lemma 09Y4: sources/stacks-etale-cohomology/pinned-original-source.tex, source lines 11745--11745. Navigation pointer invoked at proof line 11745; the counted span contains the pointer only and the cited result is not reproduced inside it, so the pointer is registered as a non-load-bearing navigation stub in the established form for Stacks material. No bibliography citation is involved.\par\medskip
\begin{lemma}
\label{lemma-constructible-constructible}
Let $X$ be a quasi-compact and quasi-separated scheme.
Let $\mathcal{F}$ be a sheaf of sets, abelian groups,
$\Lambda$-modules (with $\Lambda$ Noetherian) on $X_\etale$.
If there exist constructible locally closed subschemes $T_i \subset X$
such that (a) $X = \bigcup T_j$ and (b) $\mathcal{F}|_{T_j}$ is
constructible, then $\mathcal{F}$ is constructible.
\end{lemma}

\begin{proof}
First, we can assume the covering is finite as $X$ is quasi-compact
in the spectral topology
(Topology, Lemma \ref{topology-lemma-constructible-hausdorff-quasi-compact}
and
Properties, Lemma
\ref{properties-lemma-quasi-compact-quasi-separated-spectral}).
Observe that each $T_i$ is a quasi-compact and quasi-separated
scheme in its own right (because it is constructible in $X$; details omitted).
Thus we can find a finite partition $T_i = \coprod T_{i, j}$ into
locally closed constructible parts of $T_i$ such that
$\mathcal{F}|_{T_{i, j}}$ is finite locally constant
(Lemma \ref{lemma-constructible-quasi-compact-quasi-separated}).
By Topology, Lemma \ref{topology-lemma-constructible-in-constructible}
we see that $T_{i, j}$ is a constructible locally closed subscheme of $X$.
Then we can apply Topology, Lemma
\ref{topology-lemma-constructible-partition-refined-by-stratification}
to $X = \bigcup T_{i, j}$ to find the desired partition of $X$.
\end{proof}

\end{document}
